% Part: methods
% Chapter: proofs
% Section: starting-proofs

\documentclass[../../../include/open-logic-section]{subfiles}

\begin{document}

\olfileid{mth}{prf}{str}

\olsection{ثبوت شروع کرنا}

پر شروع کتھوں کریے؟

تہانوں کوئی گل ثابت کرن لئی دتی گئی اے، سو ثبوت وچ سبھ توں آخر اوسے گل دا ذکر ہونا چاہیدا اے (بے شک تسی شروع وچ \emph{دس سکدے او} کہ تسی ایہ گل ثابت کرو گے، پر اینوں منّی ہوئی شرط وجوں ورتنا ٹھیک نہیں)۔ جو کجھ تسی ثابت کرنا چاہندے او، اوہ اک نویں ورق دے تھلّے لکھ لو۔ ایویں تہاڈا مقصد تہاڈی نظر وچ رہے گا۔

اگے، ہو سکدا اے تہاڈے کول کجھ منّن والیاں شرطاں ہون جیہڑیاں تسی ورت سکدے او (جدوں اسیں اگلے حصے وچ تہاڈے ثبوت دی \emph{قسم} دی گل کراں گے تاں ایہ ہور صاف ہو جائے گا)۔ ایہناں نوں ورق دے اُتے لکھو تے صاف دسو کہ ایہ منّیاں ہوئیاں شرطاں نیں (یعنی جے تسی $p$ منّ رہے او، تاں «$p$ منّ لو»، یا «فرض کرو کہ $p$» لکھو)۔ آخر، ہو سکدا اے سوال وچ کجھ تعریفاں ہون جیہڑیاں تہانوں جانن دی لوڑ اے۔ شاید تہانوں اک خاص تعریف ورتن لئی آکھیا جائے، یا تہاڈیاں منّیاں شرطاں یا لوڑیندے نتیجے وچ کئی تعریفاں ہون۔ \emph{ایہناں نوں لکھ لو تے پکا کر لو کہ تہانوں ایہناں دا مطلب سمجھ آ گیا اے۔}

تسی ثبوت دی بناوٹ کیویں رکھدے او، ایہ سوال دی صورت اُتے وی ٹکدا اے۔ اگلا حصہ دسدا اے کہ جملے دی قسم مطابق ثبوت دی بناوٹ کیویں رکھنی اے۔

\end{document}
